\section{赋役、货币与边郡：把扩张的成本放回社会}
\label{sec:western-han-fiscal-society-frontier}

帝国扩张以前，编户与货币已经是朝廷调节社会的两条不同路径。前一七八年文帝下诏将田租减为三十税一，意在使战后人口、土地和生产较易重新接入户籍；这是一项可核对的诏令，不是一份全国实收账。各地的附加役使、实际税率以及未被登记的人户，仍须有地方文书才能逐项回答。轻徭薄赋因而首先是一种中央对编户恢复的政治语言，而不能直接等同每个农家的负担。

\subsection{放宽铸钱与重建控制}

前一七七年前后，文帝弛山泽禁并允许民间铸钱。中央财政有限、又希望松动秦以来的管制，诸侯和私人的铸钱空间随之扩大；钱币质量与地方财力的问题也更显著。范铸遗存能说明铸币的物质条件，不能据一处遗存推知全国市场的交易规模。田租与铸钱同被归入“宽政”，不表示它们作用在同一层面：前者面对的是可登记的土地与农户，后者则让矿产、技术、商人和诸侯国财力进入更分散的流通网络。

武帝远征和新郡经营扩大了这两类调度的压力。前一一九年收铸五铢钱、强化盐铁专营，前一一八年前后推行均输、平准，目的在于将贡赋、运输、价格与战费纳入中枢可掌握的网络。制度的先后和大体方向可以由《史记》〈平准书〉、《汉书》〈食货志〉及钱范、钱币窖藏互相校验；实施地域、官营收益和民间反应却不能被整齐地推出。均输、平准取得的是关于货物与差价的行政视野，同时也增加了地方转运和官商竞争的成本。

\subsection{轮台之后，财政争论没有终止}

前九八年前后持续整顿币制，五铢钱成为主要法定货币；出土钱币数量可以提示流通的物质痕迹，却不等于各地交易量或政策执行率。前八九年轮台诏拒绝继续在西域大规模屯田远征，显示连年屯戍、徭役和道路供给已使朝廷需要重新估量承受力。它不是边疆一概放弃的宣言，更不是此前的财政集中突然被撤销。

前八一年盐铁会议正是在这个遗产上讨论专营、民生和边防。会议没有整体废除官营，说明争论的焦点不是抽象的“国家”与“民间”二择一，而是由谁承担边郡、军队、运输和物价波动的费用。《盐铁论》保存了系统的论辩，也经过编纂，其发言不能逐字当作会议速记；但与《汉书》对读，足以看出武帝时期留下的财政工具仍在昭帝朝被重新评估。

\subsection{羌地的屯田与分类的限度}

前六一年至前六〇年赵充国主张屯田、招抚和分化羌部，反对不计成本的全面进攻。凉州与羌地的驻军要取得粮食、道路和地方合作，屯田因此既是军事方案，也是把劳力、土地和边郡行政重新接合的尝试。史书中的“羌”是官府在战争中使用的分类，不能被视为内部同质、只有一个政治目标的社会；屯田的成效和部众的具体处境也不能从将领传记直接量化。西汉的边疆财政由此显示出一个持续限制：中枢可以改变租税、钱币和转运的规则，却仍须经过地方人群、土地和交通，才能使这些规则成为可用的资源。

\subsection{材料的尺度}

本节以《汉书》〈文帝纪〉、〈食货志〉、〈西域传〉和赵充国传为年序与制度主干，并与《史记》〈平准书〉、《盐铁论》、钱范和钱币窖藏相互校验。诏令能证明政策语言，论辩能显示可说的分歧，钱币和冶铸遗存则只照见特定地点的物质活动；三者都不能单独给出基层家庭的完整收支。本节据此追踪赋役、货币与边防如何彼此牵动，而不将制度名目、出土数量或族名误作全国社会的统计。
